\documentclass[11pt,a4paper]{article}
\usepackage[T1]{fontenc}
\usepackage{isabelle,isabellesym}
\usepackage{amssymb}
\usepackage[english]{babel}
\usepackage[hidelinks]{hyperref}
\usepackage[a4paper,margin=25mm]{geometry}
\isabellestyle{it}
\hypersetup{
  pdftitle={Algebraic Invariants of the Circle, Torus, and Punctured Plane},
  pdfauthor={Arthur Freitas Ramos; David Barros Hulak; Ruy J. G. B. de Queiroz}
}
\begin{document}
\title{Algebraic Invariants of the Circle, Torus, and Punctured Plane}
\author{Arthur Freitas Ramos \and David Barros Hulak \and Ruy J. G. B. de Queiroz}
\maketitle
\begin{abstract}
  We formalize real path lifts of based loops in the circle and prove that
  their integer endpoints classify these loops up to path homotopy.  For
  the torus, the two coordinate lifts give a classification by pairs of
  integers and an explicit fundamental-group isomorphism with
  $\mathbb{Z}\times\mathbb{Z}$. Every continuous circle self-map acts on
  first singular homology by multiplication by its winding number.
  We construct a radial deformation retraction of the
  punctured complex plane onto the circle and compute the reduced homology
  of both spaces: it is isomorphic to the additive group of integers in
  degree one and trivial in every other degree.  Neither space is
  contractible.  The development also provides a reusable interface for
  path lifting in covering spaces.
\end{abstract}
\tableofcontents
\section{Introduction}
Path lifting assigns an integer to each based loop in the circle.  This
integer classifies loops up to path homotopy and provides two coordinates
for the corresponding classification on the torus.  Radial retraction
connects the circle with the punctured complex plane and transfers the
homology computation between them.  These constructions follow standard
treatments of algebraic topology, such as Hatcher~\cite{Hatcher2002}.
We build on the AFP circle
entry~\cite{FundamentalGroupCircleAFP} and use the classical
Seifert--van Kampen entry~\cite{SeifertVanKampenAFP}.

The principal contributions are the real lift-endpoint classification of
circle loops, an explicit torus fundamental-group isomorphism with
$\mathbb{Z}\times\mathbb{Z}$, and the transfer of reduced homology to the
concrete circle and punctured plane. We also prove that the actual map
induced on first singular homology by every continuous circle self-map
is multiplication by its winding number, in any integer coordinates.
The proof represents loops as singular cycles, relates concatenation to
addition using a singular triangle, and proves that winding number
annihilates singular boundaries. A homotopy through rotations removes
the basepoint restriction on the self-map.
The generic lifting interface packages the analysis library's results.
The Seifert--van Kampen statements are re-exports of the cited AFP entry;
the degree laws specialize the homology library to the abstract one-sphere.
The concrete induced-map theorem is separate from the library constant
\texttt{Brouwer\_degree2}; no additional identification with that constant,
general Hurewicz theorem, or torus homology computation is asserted here.

The covering interface is independent of the circle development.  The
product, radial-retract, and homology theories may be imported separately;
the final theory combines their interfaces.
\input{session}
\bibliographystyle{abbrv}
\bibliography{root}
\end{document}
